\section{序論}
今回の観測窓では、生成AIを支える計算基盤と、その上で動くモデル・サービスの双方に大きな話題が現れた。最も目立ったのは、NVIDIA、OpenAI、SB Energyによるオハイオ州の大規模データセンター計画である。モデル競争の外側で、電力、土地、設備投資、長期契約まで含む「AI工場」の建設が技術競争そのものになりつつあることを印象づけた。

一方、StripeによるOpenRouter買収との報道、Qwen3.8-27Bのローカル実行・車載展開、GitHubの世界的障害など、AI利用を支える経路や開発基盤にも注目が集まった。閉じた単一サービスだけでなく、モデルルーティング、オープンウェイト、エッジ推論、開発プラットフォームといった周辺層が、実用上の競争力を左右する構図がX上でも鮮明になっている。

加えて、Claudeのテキスト・ウォーターマークをめぐる議論やSeedance 2.5の利用可能性など、生成物の来歴と新しい生成体験に関する話題も継続した。以下ではGrokが収集した8件を、記載順のまま紹介する。
